\section{Appendix}

\ifmaraanonymous
\paragraph{Anonymous note.}

This anonymous version omits identity, student, course, and supervisor details. The technical appendix retains command families, evaluation protocol, project mapping, and diagnostic notes that do not identify the author.
\fi

\noindent\textbf{MSc project mapping.}\par\smallskip

\noindent\begin{tabularx}{\linewidth}{@{}p{0.30\linewidth}>{\raggedright\arraybackslash}X@{}}
\toprule
Requirement area & MARA coverage \\
\midrule
Practical implementation & Local Web UI, CLI workflows, document indexing, route-aware DocQA runtime, evidence preview, citations, and benchmark artifacts. \\
Analytical evaluation & Full-system benchmark synthesis over 41 complete artifact sets and 3,540 prediction records, with route metrics, paired deltas, bootstrap CIs, route decisions, backend metadata, and diagnostics. \\
Innovation & Controller-centred extension of a basic local QA framework using Self-RAG-inspired route decisions, evidence checks, workflow-cost traces, route switches, and abstention. \\
Ethics & Public, synthetic, or self-provided non-sensitive documents; no human participants; local-first execution and visible backend metadata. \\
\ifmaramsc
Project context & COMP702 -- M.Sc. Project (2025/26), School of Computer Science and Informatics, University of Liverpool. \\
\else
Project context & MSc delivery mapping retained in the appendix for completeness; venue-facing versions omit course identifiers, student ID, and supervisor details. \\
\fi
\bottomrule
\end{tabularx}

\ifmaramsc
For MSc delivery, these constraints define a clear prototype boundary: a usable local workbench, a controller with inspectable routes, and benchmark evidence that explains the system's behaviour. For a later submission, the strongest path is to keep the system/demo framing, include a real UI screenshot and short video, and avoid overstating results as leaderboard performance.
\else
For venue-facing versions, this mapping is supplementary project context rather than a main-paper claim. The main paper is framed as a system demo with a public release, screenshot, command surface, and route-level evaluation.
\fi

\paragraph{Full-system benchmark coverage.}

The main paper reports the three datasets that best support the demo narrative: QASPER for text-heavy scientific QA, ALCE-ASQA for citation-grounded QA, and MMDocRAG for multimodal QA. The full-system benchmark claim in the abstract is broader than this selected table. Table~\ref{tab:appendix-coverage} accounts for all 41 complete artifact sets and all 3,540 per-example prediction records in the synthesis bundle. Records count route-level prediction rows, not unique questions. Table~\ref{tab:appendix-diagnostic-results} reports the remaining diagnostic dataset results on the same 0--100 scale as the main table; the source CSV files keep the original 0--1 scores.

\begin{center}
\small
\setlength{\tabcolsep}{3pt}
\begin{tabularx}{\linewidth}{@{}lrr>{\raggedright\arraybackslash}X@{}}
\toprule
\rowcolor{gray!10}
Dataset & Sets & Records & Paper role \\
\midrule
ALCE-ASQA & 4 & 600 & Main citation-grounded dataset. \\
FinanceBench & 3 & 600 & Diagnostic numeric and financial QA dataset. \\
MMDocRAG & 20 & 600 & Main multimodal and visual dataset. \\
QASPER & 4 & 600 & Main text-heavy scientific QA dataset. \\
RAGTruth & 6 & 900 & Diagnostic groundedness and abstention dataset. \\
SlideVQA & 4 & 240 & Diagnostic visual slide QA dataset. \\
\bottomrule
\end{tabularx}
\refstepcounter{table}
\label{tab:appendix-coverage}
\smallskip
Table~\thetable: Full-system benchmark coverage.
\end{center}

\begin{table*}[t]
\centering
\small
\setlength{\tabcolsep}{4pt}
\begin{tabular}{llrrrrrr}
\toprule
\rowcolor{gray!10}
Dataset & Route & N & F1 & Native & Citation & Page hit & False abs. \\
\midrule
FinanceBench & controller\_auto & 150 & 15.10 & 8.67 & 51.00 & 60.00 & 2.67 \\
FinanceBench & crag\_guarded & 150 & 15.33 & 8.67 & 51.00 & 60.00 & 2.67 \\
FinanceBench & direct\_answer & 150 & 17.34 & 2.67 & 0.00 & 0.00 & 0.00 \\
FinanceBench & text\_rag & 150 & 14.83 & 8.67 & 48.56 & 60.00 & 7.33 \\
\midrule
RAGTruth & controller\_auto & 300 & 21.68 & 54.00 & 55.67 & -- & 0.00 \\
RAGTruth & crag\_guarded & 300 & 21.45 & 54.00 & 55.67 & -- & 1.00 \\
RAGTruth & text\_rag & 300 & 23.06 & 54.00 & 98.33 & -- & 0.00 \\
\midrule
SlideVQA & controller\_auto & 120 & 43.25 & 43.25 & 100.00 & 100.00 & 0.00 \\
SlideVQA & page\_image\_rag\_vlm & 120 & 43.25 & 43.25 & 100.00 & 100.00 & 0.00 \\
\bottomrule
\end{tabular}
\caption{Diagnostic dataset route results.}
\label{tab:appendix-diagnostic-results}
\end{table*}

\paragraph{Command and skill surface.}

MARA separates high-permission product commands from the specialist document-QA workflow. The paper's system claims refer to the shared runtime and persisted state behind these commands, not to independent QA implementations. This reduces Web/CLI drift, but contract tests remain necessary because the codebase still has entrypoints across \texttt{ktem}, \texttt{kotaemon}, and \texttt{slide\_cli}.

\begin{small}
\begin{description}[leftmargin=0pt,labelindent=0pt,style=nextline,itemsep=2pt,topsep=2pt]
    \item[App lifecycle] \texttt{MARA app init}, \texttt{MARA app doctor}, \texttt{MARA app run}.
    \item[DocQA indexing] \texttt{MARA docqa index ./docs/report.pdf}, \texttt{MARA docqa files}, \texttt{MARA docqa delete \textless file-id\textgreater}.
    \item[DocQA answering] \texttt{MARA docqa ask -{}-file report.pdf -{}-reasoning mara -{}-prompt "..."}; \texttt{MARA docqa chat -{}-file report.pdf}.
    \item[Sessions and sources] \texttt{MARA docqa sessions}, \texttt{MARA docqa resume \textless conversation-id\textgreater}, \texttt{MARA docqa sources select \textless conversation-id\textgreater{} -{}-file paper.pdf}.
    \item[Notes and artifacts] \texttt{MARA docqa notes save-answer \textless conversation-id\textgreater}; \texttt{MARA docqa artifacts generate \textless conversation-id\textgreater{} -{}-type quiz}.
    \item[Agent support assets] \texttt{MARA platform install -{}-platform codex -{}-mode full -{}-yes}; \texttt{MARA platform install -{}-platform claude-code -{}-mode full -{}-yes}; \texttt{MARA platform validate}.
\end{description}
\end{small}

{\raggedright
Focused skill wrappers include \texttt{MARA-docqa-index}, \texttt{MARA-docqa-ask}, \texttt{MARA-docqa-chat}, \texttt{MARA-docqa-resume}, \texttt{MARA-app-run}, \texttt{MARA-model-run}, and platform installation/status skills. They provide agent-facing documentation around the same CLI surface. The local Web app default is port 7860; \texttt{MARA app run -{}-port}, \texttt{GRADIO\_SERVER\_PORT}, or platform \texttt{PORT} can override it.
\par}

\paragraph{Kotaemon increment boundary.}

MARA is a branded fork of Kotaemon. The inherited substrate provides upload, model configuration, indexing, and ordinary local document QA. MARA's paper contribution is the controller-centred extension: explicit route decisions, evidence and verification traces, multimodal route interfaces, DocQA CLI reuse, platform command wrappers, and route-level benchmark artifacts. Internal package names remain compatible with the fork history, but the user-facing product surface is \texttt{MARA} / \texttt{MARA-cli}.

\paragraph{Additional route interpretation.}

Routes outside the main top-five are still part of the implemented system.
\texttt{direct\_answer} is a baseline and a useful non-document control.
\texttt{element\_rag} exposes table, formula, and element metadata but showed weak MMDocRAG quality in the full-system benchmark, consistent with sparse element evidence and high false abstention.
Graph routes are scoped routes for summaries, comparison, and global context; they are visible in controller decisions but need stronger standalone graph evaluation before being promoted to a headline route.
\texttt{visual\_retriever\_only} is best understood as a retrieval diagnostic rather than a full QA route.

\paragraph{Post-benchmark diagnostic.}

After the full-system benchmark, a targeted RAGTruth prompt-budget rerun inspected a prompt-length failure mode. The repair was verified on a small diagnostic subset, but it was not part of the same full-system protocol. The paper therefore treats it as a resolved engineering diagnostic, not as a replacement for the main RAGTruth or route-ranking results.

\paragraph{Reproducibility notes.}

The software artifact for this paper is the public MARA release \texttt{v0.0.40}\footnote{\url{https://github.com/262412/MARA/releases/tag/v0.0.40}}. The release tag identifies the application version used for the demo. Benchmark evidence is packaged separately as the companion supplementary bundle\footnote{\path{mara-full-system-benchmark-synthesis-v0.0.40.zip}}, rather than referenced through environment-specific runtime locations. In an external submission, that bundle accompanies the paper and contains a README, the synthesis report, aggregate result tables, route-decision summaries, failure/latency/backend summaries, paired-delta tables, bootstrap confidence intervals, controller-oracle regret diagnostics, win/loss/tie diagnostics, the Slurm submission ledger, the synthesis input manifest, per-example metric records, and the synthesis script. Exact filenames are listed in the bundle README.

Inside the supplementary files, runtime roots are replaced by a symbolic benchmark-run root while dataset, route, backend, score, and manifest fields are preserved. The artifact files themselves should be treated as the authority for scores and backend metadata; the external-facing paper uses the public release tag and supplement filename as its reproducibility anchors.

The run used local model and retrieval backends recorded in the benchmark metadata, including Qwen/Qwen3-8B for text generation and BAAI/bge-m3 for text retrieval. Visual route metadata must be read per run: the default local visual retriever is \texttt{local\_late\_interaction} with backend type \texttt{deterministic\_smoke}; ColPali/ColQwen-style retrievers are configured HTTP backends when available; and local Qwen3-VL generation is used only when the VLM endpoint is configured and healthy. Paper-grade external evaluators were not configured for the main result table; results are reported as local dataset-native scores and MARA diagnostic metrics.
